\vspace{-1em}
\section{Background}

\subsection{Hyperdimensional Computing}

Hyperdimensional Computing (HDC) is an emerging computational paradigm inspired by the brain’s use of distributed, high-dimensional representations for cognition. Instead of operating on low-dimensional numerical values as in conventional machine learning, HDC represents data using high-dimensional vectors, called hypervectors, often consisting of thousands of dimensions. Information is encoded and manipulated through simple, mathematically well-defined operations on these hypervectors, enabling highly efficient and robust computation. HDC systems are naturally error-tolerant due to their distributed structure and can learn from limited data without the need for expensive, iterative optimization methods like backpropagation used in deep learning. This makes HDC particularly suitable for resource-constrained environments and applications requiring fast, on-device learning or inference. While traditional deep learning models often require large amounts of data, substantial compute resources, and energy-intensive training pipelines, HDC offers a lightweight but effective alternative with lower computational overhead, faster training, and inherent suitability for hardware acceleration. These properties make HDC a compelling approach for applications in edge AI, biomedical signal processing, and other domains where efficiency, robustness, and interpretability are critical. 

The increased performance and inherent error resilience of HDC come at a cost: achieving the desired accuracy for a given application typically requires extensive exploration of hyperparameters and algorithmic strategies to compensate for the simplified learning mechanisms employed by HDC.
A typical HDC application consists of three main stages: encoding, training, and inference. In the encoding stage, raw input data is transformed into high-dimensional hypervectors, capturing the essential features of the data in a distributed representation. During training, hypervectors corresponding to inputs from the same class are combined to form class hypervectors or prototype representations, typically using simple aggregation operations such as vector addition and normalization. Finally, in the inference stage, a query hypervector is compared to the stored class hypervectors using a similarity or dissimilarity metric, commonly cosine similarity or Hamming distance, to identify the closest match and produce the predicted label. Each of these stages is highly parallelizable, allowing HDC applications to efficiently leverage modern heterogeneous hardware, and the distributed nature of hypervectors provides robustness to noise and approximation throughout the pipeline.

\subsection{HDC++}
\label{bg:hdc++}

HDC++~\cite{hpvm-hdc} is a domain-specific, high-level programming language for hyperdimensional computing, embedded within C++. It exposes core HDC operations as first-class primitives, enabling productive and expressive development of HDC applications. Computations in HDC++ operate on hypervectors, representing individual encoded features or class data, and on hypermatrices, representing collections of hypervectors. Unlike prior programming approaches for HDC, HDC++ supports retargetable compilation across a wide range of hardware, including CPUs, GPUs (with optimized libraries), and custom HDC accelerators, using the HPVM-HDC compiler, a retargetable compiler developed specifically for heterogeneous HDC workloads. In terms of development effort, HDC++ applications typically require about twice the lines of code compared to an equivalent Python implementation targeting CPUs, but are significantly more concise than equivalent CUDA implementations. Crucially, a single HDC++ application can be compiled for multiple hardware targets, including accelerators for which no prior implementation exists, eliminating the need for maintaining multiple target-specific versions of the same application.

Another key feature of HDC++ is the ability for developers to control software approximation through concise source code annotations, typically requiring only one or two lines per HDC primitive. These annotations instruct the HPVM-HDC compiler to generate hardware-specific code with the specified approximation transformations applied. The details of these transformations are described in Section~\ref{sec:hdc-approximations}. This compact annotation mechanism enables productive exploration of the approximation space inherent in HDC, allowing developers to achieve the desired accuracy manually while maintaining control over performance and efficiency trade-offs. Listing~\ref{lst:hdcpp} describes a simple HDC++ example which encodes input data using random projection encoding, calculates the dissimilarity using Hamming distance and then identifies the most dissimilar label using argument minimum. The approximation annotation on line 12 instructs the HPVM-HDC compiler to generate target-specific code which calculates Hamming distance using only the first 1024 elements of data.

% \lstinputlisting[label={lst:hdcpp}, caption={HDC++ example for random-projection encoding and dissimilarity between query and class hypervectors.}]{code_snippets/hdcpp_snippet.tex}

\lstinputlisting[label={lst:hdcpp}, caption={HDC++ example for random-projection encoding and dissimilarity between query and class hypervectors.}]{code_snippets/other_snippet.tex}
        

\subsection{HPVM-HDC}
\label{bg:hpvm-hdc}
HPVM-HDC~\cite{hpvm-hdc} is a retargetable compiler infrastructure designed specifically for hyperdimensional computing (HDC) applications. Built on top of the HPVM~\cite{kotsifakou2018hpvm, ejjeh2022hpvm} heterogeneous compiler framework, HPVM-HDC provides an intermediate representation (IR) that captures the high-level semantics of HDC operations, including encoding, training, and similarity-based inference. The HPVM-HDC IR abstracts away hardware-specific details while preserving the structure of HDC primitives, enabling the compiler to perform optimizations such as parallelization and target-specific code generation. 
Each HDC++ language primitive corresponds to a unique HPVM-HDC IR operation, allowing the intermediate representation to retain the domain-specific semantics of HDC within the compiler. Internally, the HPVM-HDC compiler translates HDC++ applications into the HPVM IR, which provides constructs for representing multiple forms of parallelism, including task-level, data-level, and streaming parallelism. Parallel programs are represented as a hierarchical directed acyclic graph (DAG), where nodes fall into two categories: internal nodes, which capture hierarchical parallelism by containing entire subgraphs, and leaf nodes, which represent individual units of computation expressed in LLVM IR. Edges between nodes denote logical data transfers, meaning explicit copies may or may not be required. Additionally, each node can be annotated with one or more hardware targets, enabling the compiler to generate code for multiple backends and fully exploit heterogeneous platforms such as CPUs, GPUs, and (Intel) FPGAs.

In addition to generating HPVM IR for retargetable heterogeneous compilation, HPVM-HDC implements two parameterized software approximation transformations at the IR level. The first, Automatic Binarization Propagation, performs a quantization-like transformation on hypervectors, representing the resulting data as single-bit packed vectors and inter-procedurally updates computations and allocations as required. The second, Reduction Perforation, applies parameterized loop perforation by regularly skipping iterations in HDC operations that involve reductions. Both of these software approximations are described in detail in Section~\ref{sec:hdc-approximations}. HPVM-HDC also supports compilation for specialized HDC accelerators, including an HD-ASIC and an HD Resistive RAM (ReRAM) accelerator. However, it does not currently exploit hardware-specific approximation strategies for these accelerators, a capability we introduce in this work.

\subsection{ReRAM and SpecPCM}

Resistive RAM (ReRAM) and phase-change memory (PCM) are two classes of non-volatile memory that support in-memory computing, the execution of arithmetic directly within the memory array without moving data to a separate processing unit. Both technologies store information as analog resistance states in dense crossbar arrays, and both can perform vector operations such as dot products and Hamming distance through the physics of current summation along array columns. These properties make them a natural substrate for HDC, where the core computational patterns reduce to parallel element-wise and reduction operations over long vectors.

\textbf{ReRAM-based HDC acceleration.} In a ReRAM crossbar, each memristive device is programmed to a conductance proportional to a hypervector element. A read operation applies input voltages to the rows and senses the summed column currents, effectively computing a dot product in a single analog step. FSL-HD \cite{ReRAMAcc} exploits this structure as a few-shot learning accelerator that performs both tensorized encoding and Hamming-distance inference entirely within the ReRAM array. By keeping hypervectors resident in memory and avoiding off-chip data movement, FSL-HD achieves substantial speedups and energy savings over CPU and GPU baselines, while also improving on the area efficiency of the earlier SAPIENS design \cite{SAPIENS}. The key insight from an approximation standpoint is that ReRAM reads are inherently imprecise: device-to-device variability, write noise, and limited ADC resolution all introduce errors, yet HDC accuracy is largely preserved because the final prediction depends on aggregate statistics across thousands of dimensions rather than on any single element.

\textbf{PCM-based HDC acceleration.} SpecPCM \cite{SpecPCM2024} takes a complementary approach using phase-change memory. The architecture offloads distance computations, spectral clustering distances and Hamming similarity for database search into PCM crossbar arrays organized as parallel memory banks with a 2T2R cell structure. SpecPCM integrates two distinct superlattice PCM device types within the same accelerator: one optimized for fast, low-voltage programming suited to clustering workloads, and another providing long retention and low read noise for reliable similarity search. SpecPCM further increases effective storage density through multi-level cells (MLC), packing multiple hypervector dimensions into a single device using a dimension-packing scheme. Stages not suitable for analog execution (e.g., argument selection) are handled by near-memory ASIC logic, enabling significant speedups and energy-efficiency gains over software baselines.

ReRAM and PCM accelerators provide tunable parameters (e.g., ADC resolution, write-verify depth, MLC levels, quantization scale) that trade computation accuracy for latency, energy, and area by introducing controlled noise. Due to HDC’s robustness to noise, this has minimal impact on accuracy. In this work, \pname{} includes these hardware knobs in its approximation search space alongside software techniques to jointly optimize accuracy, performance, and energy.
 
% Resistive RAM (ReRAM) and phase-change memory (PCM) technologies enable in-memory computing accelerators for hyperdimensional computing (HDC) by performing key operations directly within dense memory arrays. A ReRAM-based HDC accelerator stores hypervectors in memristive crossbar arrays and uses analog conductance summation to accelerate operations such as cosine or Hamming similarity. For example, Xu et al.~\cite{FSL-HD} describe an HDC accelerator built on large ReRAM crossbars that supports ``tensorized'' encoding (a variant of random projection encoding) and summation-based one-shot training, with Hamming-distance inference executed directly in analog memory. This design performs both encoding and inference within the ReRAM array, leveraging the technology's high on-chip density and fast read times to achieve substantial parallelism. Their FSL-HD prototype for few-shot learning on ReRAM reports significant speedups and large improvements in energy efficiency over CPU and GPU baselines, even outperforming a prior ReRAM-HDC design (SAPIENS~\cite{SAPIENS}) while using considerably less area. Overall, ReRAM-based analog in-memory computing is inherently well aligned with the binary, element-wise computational structure of HDC workloads.

% By contrast, SpecPCM is a PCM-based in-memory accelerator that pushes HDC computing into analog phase-change memory devices. The SpecPCM architecture offloads distance computations to the memory. For instance, spectral-clustering distances and Hamming similarity for database search are executed inside PCM crossbar arrays, each organized as a large memory bank with a standard 2T2R cell structure. Multiple banks operate in parallel to exploit data-level parallelism. Fan et al.~\cite{SpecPCM2024} integrate two distinct superlattice PCM device types, each tailored to a specific HDC operation. One device type enables fast, low-voltage programming for clustering workloads, while the other provides long retention and low read noise required for reliable similarity search. The accelerator supports multi-level cells (MLC) to increase effective density, combined with a dimension-packing technique that maps multiple hypervector dimensions onto a single cell. The design integrates custom near-memory ASIC blocks for stages that are not amenable to analog computing, while the memory arrays handle the bulk of HDC distance calculations. As a result, SpecPCM achieves significant gains over conventional processors, offering substantial speedups for clustering and search, along with major improvements in energy efficiency relative to state-of-the-art CPU and GPU implementations. Analog in-memory acceleration is highly effective when paired with error-tolerant algorithms such as HDC.
